\thispagestyle{empty}
%\input{00-Border/border}

\vspace{-2cm}
\begin{center}
\textbf{\large \textcolor{BrickRed}{چکیده} }
\end{center}
%\vspace{-1.5cm}

منابع انرژی الکتریکی همانند باتری‌‌‌های با پایه سرب-اسید، نیکل، لیتیوم و پیل‌‌‌های سوختی به‌‌‌طور گسترده‌‌‌ای توسط ابزارها و ماشین‌‌‌های الکتریکی مورد استفاده قرار می‌‌‌گیرند. افزایش روزافزون قیمت نفت و اتمام منابع انرژی‌‌‌های فسیلی، توجه تمامی کارشناسان را به‌‌‌سمت استفاده از منابع تولید انرژی غیرفسیلی، خصوصاً باتری‌‌‌ها و پیل‌‌‌های سوختی، جلب کرده و جایگزین کردن آنها در خودروها و دیگر کاربردها را از لحاظ اقتصادی توجیه کرده‌‌‌است. به همین دلیل شناخت هرچه کامل‌تر رفتار باتری‌ها و حل مشکلات آنها ضروری بنظر می‌رسد. مشکلات باتری‌ها را از چند جهت می‌توان ارزیابی نمود که عبارتند از: محتوای انرژی، توان، طول عمر، بروز پدیده گریز حرارتی و پدیده چینه‌‌‌بندی الکترولیت در باتری‌‌‌های با الکترولیت شناور. 
این مشکلات که اغلب در هنگام استفاده از باتری‌ها تحت سیکل‌های شارژ و دشارژ بالا رخ می‌دهد، زمینه‌هایی است که محققان برای برطرف نمودن آنها در تلاش هستند.
در این گزارش که تحت عنوان
\textbf{شبیه‌سازی باتری سرب-اسید به روش رتبه کاسته}
است تلاش می‌شود تا با استفاده از روش‌های مدل‌سازی مختلف مدل‌هایی از باتری ارائه شود که بوسیله آنها بتوان رفتار باتری را تحت شرایط مختلف پیشگویی کرد.  از آنجائیکه این معادلات بسیار غیر خطی هستند حل آنها توسط روش‌های رایج دینامیک سیالات محاسباتی مانند روش حجم محدود و روش اختلاف محدود بسیار زمان‌بر است
%اگرچه با پیشرفت تکنولوژی پردازشگرهای رایانه‌ای و استفاده از روش‌های موازی و یا استفاده از پردازشگرهای گرافیکی تا حد زیادی به سرعت پردازش اطلاعات افزوده است، اما باز هم بدلیل پیچیدگی رفتار و معادلات حاکم بر باتری سرب-اسید این روش‌ها رضایت بخش نبوده
 و لذا نیاز به روش‌های عددی دیگری است که بتوان با دقت بسیار بالا، زمان محاسبات را تا حد زیادی کاهش داد.
در این راستا، یکی از بهترین روش‌ها، روش رتبه کاسته%
 \LRfootnote{Reduced Order Modeling} 
است. روش‌های رتبه کاسته به‌منظور کاهش تعداد درجات آزادی مسائل غیرخطی به‌کار می‌روند.
%در اینگونه مسائل، معادلات حاکم شامل یکسری معادلات دیفرانسیلی جزئی غیرخطی هستند.
روش‌های رتبه کاسته را به شیوه‌های مختلفی می‌توان انجام داد که از جمله آنها می‌توان به مدل‌سازی رتبه‌کاسته با استفاده از مدهای ویژه جریان و مدل‌سازی رتبه کاسته با استفاده از توابع متعامد سره \lr{POD}%
 \LRfootnote{Proper Orthogonal Decomposition} 
 اشاره نمود.

  برای نیل به این منظور در این رساله ابتدا معادلات حاکم بر دینامیک باتری سرب-اسید بدست آورده شده که شامل معادلات ناویر استوکس، معادله انرژی و معادلات الکتروشیمیایی هستند. سپس این دستگاه معادلات توسط روش‌های رتبه‌کاسته و بر اساس توابع \lr{POD} حل می‌شوند.
 
 توابع \lr{POD} در واقع بهترین بیان خطی از داده‌های موجود می‌باشند. این داده‌ها می‌توانند به‌صورت تجربی، تحلیلی یا عددی به‌دست آیند. در این رساله داده‌های لازم برای تشکیل توابع \lr{POD} از مدل‌های عددی به‌دست می‌آیند.
 مدل‌سازی رتبه کاسته بر اساس توابع \lr{POD} خود شاخه‌های گوناگونی دارد.
  در نوع اول این مدل‌سازی معادلات حاکم روی توابع \lr{POD} تصویر می‌شوند که در این صورت به آن تصویر گالرکین یا همان تصویر تحلیلی روی توابع می‌گویند. در اینصورت معادلات حاکم به معادلات دیفرانسیلی معمولی تبدیل می‌شوند. در حالت دیگر معادلات حاکم ابتدا گسسته شده و سپس معادلات گسسته‌سازی شده روی توابع \lr{POD} تصویر می‌شوند. در این صورت نیز بسته به نوع گسسته‌سازی، معادلات دیفرانسیلی مشتق جزئی حاکم به معادلات دیفراسیل معمولی و یا به معادلات جبری تبدیل می‌شوند که در این صورت براحتی و با سرعت بالا حل می‌شوند.

\thispagestyle{empty} 
\hspace{-0.65cm}%
\textbf{\textcolor{BrickRed}{کلمات کلیدی}} :
باتری سرب-اسید، شبیه‌سازی عددی، مدل‌سازی سازی رتبه‌کاسته، توابع متعامد سره

%\thispagestyle{empty}
%\input{00-Border/border}
\clearpage
%\pagebreak